

\chapter{Stand des Wissens / Stand der Technik} \label{cha:standDesWissens}
Die vorliegende Masterarbeit untersucht den Einfluss von Large Language Models (LLMs) wie ChatGPT auf das Lernen und Prüfen von Programmieranfänger*innen. Der Forschungsstand umfasst mehrere miteinander verknüpfte Bereiche: die besonderen Herausforderungen des Programmierlernens, den Einsatz von LLMs in der Programmierausbildung, deren Rolle in Prüfungs- und Bewertungsszenarien sowie didaktische Faktoren. In diesem Kapitel soll die einschlägige Literatur zusammengeführt werden.


\section{Herausforderungen beim Programmierenlernen im tertiären Bildungsbereich}
Das Erlernen von Programmierkompetenzen als besonders herausfordernd. Zahlreiche Studien zeigen, dass Einführungsveranstaltungen in der Informatik überdurchschnittlich hohe Dropout- und Durchfallsquoten aufweisen. \textcite[S.~6]{SystematicLiteratureReview2019} identifizieren in ihrer systematischen Literaturübersicht zentrale Einflussfaktoren, die den Erfolg in Einsteigerkursen bestimmen, darunter mathematische Kompetenzen, Problemlösefähigkeiten, Abstraktionsvermögen sowie die Fähigkeit, Fehlermeldungen zu interpretieren und systematisch zu debuggen. Weitere Arbeiten, etwa von \textcite[S.~2]{TeachingIntroductoryQuantitive} oder \textcite[S.~1]{studentPerformance}, zeigen, dass insbesondere die ersten Lernschritte also das Strukturieren eines Problems und der Übergang von natürlicher Sprache zu formalen Anweisungen sowie der Umgang mit Syntax als besonders frustrierend erlebt werden.

Auf curricularer Ebene bestehen zusätzliche Herausforderungen: Lehrende müssen mit unterschiedlichen Vorkenntnisse umgehen, sinnvolle Aufgabenformate entwickeln und auf eine angemessene Balance zwischen Theorie, Codepraxis und Feedback achten. \textcite[S.~7]{SystematicReviewFinnland} sowie \textcite[S.~7-8]{DesigningAndDeploying} zeigen, dass die Wahl der Programmiersprache und der didaktischen Methode (z. B. pair programming, projektorientiertes Lernen, tutorielle Unterstützung) entscheidend dazu beiträgt, wie gut Studierende grundlegende Konzepte wie Kontrollstrukturen, Datenstrukturen oder Objektorientierung verstehen.


\section{Die Verwendung von LLMs in der Programmierausbildung}
Mit dem Aufkommen von generativen Modellen wie GPT-4, Claude oder Gemini hat sich das Programmierlernen grundlegend verändert. LLMs können Code generieren, erklären, korrigieren und optimieren. Sie können Fehlermeldungen interpretieren, Beispiele erzeugen und Lernende durch ganze Problemlösungsprozesse begleiten. \textcite[S.~1]{DennyEra} beschreiben diesen Wandel als einen Paradigmenwechseln, der sowohl Chancen für individualisiertes Lernen als auch Herausforderungen für Didaktik und Prüfungsdesign mit sich bringt.

Eine Reihe von Studien befasst sich mit den unmittelbaren Lerneffekten. \textcite[S.~1]{RobotsAreComming} sowie \textcite[S.~1]{KazemitabaarCodeAid} zeigen, dass Studierende mit LLM-Unterstützung typischerweise schneller zu funktionierenden Lösungen gelangen und weniger Frustrationspunkte erleben – insbesondere in frühen Lernphasen. Ebenso deuten mehrere Arbeiten darauf hin, dass die Fähigkeit der Modelle, Fehlermeldungen zu erklären \textcite[S.~1]{LeinonenErrorMessages}, erheblich zur Reduktion von Debugging-Zeiten beiträgt.

Gleichzeitig wird in der Forschung intensiv diskutiert, ob die Nutzung solcher Systeme zu oberflächlichen Lernprozessen führt. \textcite[S.~1]{BeckerProgrammingHard} argumentieren, dass Programmieren möglicherweise nicht mehr so schwer ist wie früher, weil LLMs viele kognitive Hürden übernehmen. Damit stellt sich die Frage, welche Kompetenzen Studierende tatsächlich erwerben, wenn wesentliche Teile der Problemlösung ausgelagert werden.

Einen wichtigen Beitrag für die Praxis liefern zudem spezialisierte Systeme wie CodeHelp von \textcite[S.~1]{CodeHelp} oder CodeAid von \textcite[S.~3]{KazemitabaarCodeAid}, die LLMs mit didaktischen Leitplanken kombinieren. Diese Tools verhindern, dass der User direkt eine Code-Lösung erhält. Sie bieten stattdessen erklärende Unterstützung, wodurch eigenständiges Denken angeregt werden soll. Erste Studien zu CodeAid und ähnlichen Classroom-Assistenten belegen, dass LLMs in realen Kursen bereits als standardisierte Unterstützungswerkzeuge eingesetzt werden können – allerdings nur, wenn klare Regeln zur Nutzung definiert werden.



\section{Generative KI in Prüfungen und Bewertungsszenarien}
Der Einsatz von LLMs stellt traditionelle Prüfungsformate grundlegend infrage. Studien wie die von \textcite[S.~6]{PintoLLMsForEducation} zeigen, dass LLMs offene Fragen in vielen Fällen ähnlich wie menschliche Bewertende beurteilen können - allerdings mit typischen Abweichungen: Manche Modelle sind strenger, andere übersehen subtile Fehler. Damit wird die Frage nach Fairness, Reliabilität und Transparenz solcher Bewertungsmethoden zentral.

Parallel dazu diskutieren \textcite[S.~1]{DennyPromptProblems2024}, \textcite[S.~1]{MahonGuidelines} und \textcite[S.~108]{PratherRobotsAreHere2023}, welche Aufgabenformate noch geeignet sind, individuelle Kompetenzen im LLM-Zeitalter zu prüfen. Aufgaben, die primär auf reiner Codeproduktion beruhen, verlieren an Aussagekraft, da LLMs solche Lösungen oft vollständig generieren. Stärker gefragt sind daher Problemanalyse, Codebewertung, Fehlersuche, Designentscheidungen und das Verfassen erklärender Texte. Diese Komponenten gelten als „LLM-resistent“, weil sie Verständnis und Reflexion voraussetzen und nicht vollständig automatisierbar sind.


\section{LLM-unterstütztes Programmierenlernen}
Mehrere kontrollierte Studien untersuchen, wie sich der Einsatz großer Sprachmodelle auf das Programmierlernen von Anfänger:innen auswirkt. Eine häufig zitierte Arbeit stammt von \textcite[S.~11]{KazemitabaarSupportLearners2023}: In einem randomisierten Experiment mit 69 Beginner:innen zeigte sich, dass Lernende mit Zugriff auf einen Codex-basierten Code-Generator signifikant häufiger funktionsfähige Programme erzeugten und ihre Aufgaben schneller lösten als die Kontrollgruppe. Gleichzeitig blieben ihre Leistungen bei nachfolgenden manuellen Code-Modifikationen stabil, was ein Hinweis darauf ist, dass der Einsatz von LLMs nicht automatisch zu Abhängigkeit führt.

Eine weitere Analyse von \textcite[S.~1]{kazemitabaar_how_2023} liefert Aufschluss über die Interaktionsmuster von Programmier-Einsteiger*innen bei der Verwendung von OpenAIs Codex. Anhand von Log-Daten identifizierten die Autoren vier distinkte Nutzungsstrategien bei Novizen: „AI Single Prompt“ (Generierung der gesamten Lösung), „AI Step-by-Step“, „Hybrid“ und „Manual“. Die Ergebnisse deuten darauf hin, dass die Art der Interaktion entscheidend für den nachhaltigen Lernerfolg ist: Der „Hybrid“-Ansatz, der eine Mischung aus eigenständigem Schreiben und gezielter KI-Nutzung darstellt, zeigte positive Trends hinsichtlich der Ergebnisse in späteren Retention-Tests. Wohingegen der „AI Single Prompt“-Ansatz eher mit negativen Resultaten korrelierte. Dies unterstreicht, dass der didaktische Mehrwert von LLMs maßgeblich vom Grad der Selbstregulation (Self-Regulation) und der eigenen kognitiven Beteiligung abhängt \parencite[S.~12]{kazemitabaar_how_2023}.

\textcite[S.~3-6]{XueIntro2024} untersuchten den Einsatz von ChatGPT 3.5 in einem CS1-Java-Kurs. Obwohl ChatGPT syntaktische Fehler reduzieren konnte, zeigten sich keine signifikanten Vorteile in tiefergehenden konzeptuellen Aufgaben. Die Autor:innen betonen, dass LLMs ohne didaktische Einbettung nur begrenzt Lernzuwächse erzeugen.
Ein weiteres Experiment von \textcite[S.~6-8]{TobaTests2024} demonstriert vor allem Effizienzgewinne: Mit ChatGPT lösten Studierende Programmieraufgaben etwa doppelt so schnell wie ohne. Die Codequalität blieb jedoch vergleichbar. Gleichzeitig weist die Studie auf Risiken wie missbräuchliche Nutzung in Prüfungssituationen hin.

Eine systematische Meta-Analyse von \textcite[S.~12-15]{Alanazi2025} bestätigt diese Befunde: KI-gestützte Lernsettings verbessern im Durchschnitt Performance und Bearbeitungszeit deutlich, zeigen aber keine konsistenten Effekte auf tiefes Konzeptverständnis.



\section{User-Experience-Messungen von LLM-Chatbots mittels User Experience Questionnaires (UEQ)}
Die Evaluierung von LLM-basierten Chatbots erfordert Messinstrumente, die über reine Leistungsmetriken hinausgehen und die subjektive Wahrnehmung der Interaktion erfassen. Als Standardinstrument hat sich hierbei der User Experience Questionnaire (UEQ) etabliert, da er sowohl pragmatische (zielgerichtete) als auch hedonische (nicht-zielgerichtete) Qualitätsaspekte abdeckt \parencite[S.~2-3]{Schrepp2023UEQHandbook}.

Aktuelle Studien bestätigen die Eignung dieses Messinstruments für Chatbots: So konnten \textcite[S.~17]{Mendoza2022} in einer Untersuchung im schulischen Umfeld nachweisen, dass der UEQ verlässliche Aussagen über die Akzeptanz von Lernassistenten liefert, wobei insbesondere die Dimensionen der Effizienz und der Verlässlichkeit (Efficieny und Dependability) als kritische Erfolgsfaktoren identifiziert wurden. \textcite[]{Ambacher2024} hebt in seiner Analyse von LLM-User-Interfaces hervor, dass eine mangelnde Transparenz über die Funktionsweise des Modells oder unerwartetes Systemverhalten direkt die Dimenson Verlässlichkeit negativ beeinflusst. Wird das System als unberechenbar wahrgenommen, sinkt die Akzeptanz, selbst wenn die generierten Antworten technisch hilfreich sein könnten. Demgegenüber verweist \textcite[S.~6]{Shahu2025} auf das Potential von LLMs, die pragmatische Qualität (Effizienz, Verständlichkeit und Verlässlichkeit) einer Aufgabe oder eines Produkts stark zu steigern. Ihre Arbeit handelt von einem GIS-Assistenten namens "MapGPT", der durch Implementierung natürlichsprachlicher Eingaben die Barrieren zur Nutzung dieser komplexer Software senken konnte. Dieser Sachverhalt schlug sich in positiven Werten auf den UEQ-Skalen für Efficiency (Effizienz) und Perspicuity (Verständlichkeit) nieder. \textcite[S.~5-6]{Lee2024} verwendeten auch den UEQ als Messinstrument für die Evaluierung ihres Chatbots "ChatFive" und konnten in der Dimension Attraktivität sowie dem Bereich Stimulation signifikant höhere Werte erzielen.

Folglich erweist sich der UEQ als das methodisch adäquate Instrument, um nicht nur die funktionale Gebrauchstauglichkeit (Pragmatik), sondern auch die durch generative KI neu entstehenden emotionalen Interaktionsqualitäten (Hedonik) in einem erprobten Messmodell vergleichbar zu erfassen.


\section{Zusammenfassung}
TODO
